% !TeX program = lualatex
% Standalone research notebook for original catalogue problem 64.
% Compile twice: lualatex 64.tex
% Original entry retained; research subsection and [E] references added.
% No external bibliography, image, data, or style file is required.
\documentclass[11pt,a4paper]{article}
\usepackage[margin=24mm,headheight=14pt]{geometry}
\usepackage{fontspec}
\setmainfont{Latin Modern Roman}
\setsansfont{Latin Modern Sans}
\setmonofont{Latin Modern Mono}
\usepackage{amsmath,amssymb,mathtools,amsthm}
\usepackage{unicode-math}
\setmathfont{Latin Modern Math}
\usepackage{microtype}
\usepackage{enumitem,xurl,needspace}
\usepackage[dvipsnames]{xcolor}
\usepackage{hyperref}
\hypersetup{unicode=true,colorlinks=true,linkcolor=MidnightBlue,
 urlcolor=MidnightBlue,bookmarksnumbered=true,
 pdftitle={Research notebook - Problem 64},
 pdfauthor={Research attempt; not independently refereed}}
\usepackage{bookmark}
\usepackage{titlesec}
\titleformat{\section}{\Large\bfseries\raggedright}{\thesection}{0.7em}{}
\usepackage{fancyhdr}
\pagestyle{fancy}\fancyhf{}
\fancyhead[L]{\small\sffamily Research notebook}
\fancyhead[R]{\small\sffamily Problem 64}
\fancyfoot[L]{\scriptsize 27 September 2026}
\fancyfoot[C]{\thepage}
\fancyfoot[R]{\scriptsize Unrefereed research attempt}
\setlength{\parindent}{0pt}
\setlength{\parskip}{4pt plus 1pt}
\setlength{\emergencystretch}{3em}
\clubpenalty=10000
\widowpenalty=10000
\displaywidowpenalty=10000
\setcounter{secnumdepth}{1}
\setlist[itemize]{leftmargin=3em,itemsep=2pt,topsep=3pt,parsep=0pt,before=\small}
\allowdisplaybreaks
\newcommand{\N}{\mathbb{N}}
\newcommand{\Z}{\mathbb{Z}}
\newcommand{\Q}{\mathbb{Q}}
\newcommand{\R}{\mathbb{R}}
\newcommand{\C}{\mathbb{C}}
\newcommand{\F}{\mathbb{F}}
\newcommand{\A}{\mathbb{A}}
\newcommand{\PP}{\mathbb P}
\newcommand{\E}{\mathbb{E}}
\newcommand{\eps}{\varepsilon}
\newcommand{\dd}{\,\mathrm d}
\newcommand{\sref}[2]{\hyperlink{source:#1:#2}{[S#2]}}
\newcommand{\eref}[2]{\hyperlink{extra:#1:#2}{[E#2]}}
\newif\ifshowcatalogue
\showcataloguetrue
\newcommand{\cataloguescope}[1]{\ifshowcatalogue
 \paragraph{Exact scope in the supplied catalogue.}
 \begin{quote}\small #1\end{quote}\fi}
\newtheorem{notebooklemma}{Lemma}[section]
\newtheorem{notebookproposition}[notebooklemma]{Proposition}
\numberwithin{equation}{section}
\begin{document}
\setcounter{section}{63}
% ================================================================
% problemId: problem.vojtas-main-conjecture-on-height-inequalities-in-diophantine-approximation
% source releaseRank: 64; source releaseStatus: open
% exactTarget SHA-256: ca8ce85ca40e0f4edecdc8bd74fd01f951a783c39f1289ff4f19388614375957
\clearpage
\section{\texorpdfstring{Vojtas Main Conjecture on Height Inequalities in Diophantine Approximation}{Vojtas Main Conjecture on Height Inequalities in Diophantine Approximation}}\label{64}
\subsection{Definitions and mathematical statement}
Let $h_A$ and $h_{K_X}$ be logarithmic Weil heights for the indicated divisors, defined up to bounded functions. For local Weil functions $\lambda_{D,v}$, put $m_{D,S}(P)=\sum_{v\in S}\lambda_{D,v}(P)$ with normalized local weights. A simple normal-crossings divisor has components meeting locally like distinct coordinate hyperplanes. For every $\varepsilon>0$, the assertion is that some proper closed $Z\subset X$ and constant $C$ satisfy
\[
 m_{D,S}(P)+h_{K_X}(P)<\varepsilon h_A(P)+C
 \qquad(P\in X(k)\setminus Z).
\]
Enlarge $Z$ to include the support of $D$, where the proximity function is not finite. The constant and exceptional set may depend on the fixed data and $\varepsilon$, but not on $P$.
\par\smallskip\textit{Formulation and concept references:} \sref{64}{1}.
\cataloguescope{Let k be a number field, X a smooth projective variety over k, K\_X a canonical divisor, A an ample divisor, D an effective simple normal-crossings divisor, and S a finite set of places containing the archimedean places. Prove that for every epsilon > 0 there are a proper Zariski-closed subset Z of X and a constant C such that m\_\{D,S\}(P)+h\_\{K\_X\}(P) < epsilon h\_A(P)+C for every P in X(k) outside Z.}
\subsection{Short English statement}
Can unusually close rational approximation to a divisor be uniformly bounded by Vojta's height inequality outside a proper exceptional set?
% BEGIN RESEARCH INSERTION
\subsection{Research attempt}

\paragraph{Current frontier.}
Yasufuku's Conjecture 8 is the fixed-number-field height inequality recorded above; the paper's gcd applications do not prove the universal assertion \sref{64}{1}. Chen--Garcia-Fritz--Mathur--Pasten's 2026 degeneration manuscript treats specified integral-point finiteness and degeneracy problems, while Abe's recent manuscript treats a function-field analogue \eref{64}{2}, \eref{64}{3}. Neither replaces an inequality on all rational points outside one exceptional set. We will use generalized Chevalley--Weil, but not assume a bounded-degree Vojta inequality \eref{64}{1}, Theorem 23.15.

\paragraph{Attack route.}
A Subspace-Theorem approach needs enough suitably positioned divisors, which an arbitrary variety need not provide. Passing to a finite cover could improve the geometry, but lifts of rational points may run through infinitely many fields. The route investigated here is descent from a useful cover. We first isolate the benign étale case, then compute how ramification compensates for the discriminant of the lift. The calculation identifies precisely which uniform height theorem would still be necessary upstairs.

\paragraph{Conventions.}
Use absolute logarithmic heights and normalized local sums, invariant under extending the field of definition. For $L/k$ finite, put
\[
 d_k^{\mathrm{abs}}(L)=\frac{\log|D_L|}{[L:\Q]}-
                          \frac{\log|D_k|}{[k:\Q]},
 \qquad d_k^{\mathrm{abs}}(Q)=d_k^{\mathrm{abs}}(k(Q)).
\]
For an effective divisor $R$, write $N_{R,S_L}(Q)=h_R(Q)-m_{R,S_L}(Q)$ up to bounded functions, with $S_L$ all places over $S$. Bounded errors below are uniform on the indicated bounded-degree points. The normal-crossings boundaries in this argument are reduced, as in the source convention; arbitrary multiplicities are not appended to the statement.

\paragraph{Editorial note on boundary multiplicities.}
The reduced convention matters: Vojta's Definition 14.1 explicitly includes it and notes that terminology varies \eref{64}{1}. A support-only interpretation allowing arbitrary multiplicity would give a different, false assertion. On $\PP^1/\Q$, take $D=3[0]$, $K_X=-2[\infty]$, $A=[\infty]$, $S=\{\infty\}$, and $P=1/m$. Then $m_{D,S}(P)+h_{K_X}(P)=\log m+O(1)$ while $h_A(P)=\log m$. For $\eps<1$, no constant and proper closed subset can handle all $m$, since such a subset of $\PP^1$ is finite. This is an interpretation warning, not a counterexample to the reduced source formulation.

\paragraph{Attempt I: étale descent.}
Let $f:Y\to X$ be finite surjective étale of degree $d$, with both varieties smooth projective and geometrically integral. Set $D_Y=f^*D$ and $A_Y=f^*A$. Then $D_Y$ is SNC and $K_Y=f^*K_X$. Only finitely many number fields are needed to define all lifts of points of $X(k)$. To justify the uniformity, spread the projective varieties and the finite étale map outside a fixed finite set of primes. Properness extends each rational point over that ring of integers; its fiber is unramified outside the fixed set and has degree at most $d$. Bounded local degrees bound the discriminant exponents at the remaining primes, and Hermite--Minkowski gives the finite list. This is the projective Chevalley--Weil argument \eref{64}{1}, Sections 12 and 23.

Suppose the original rational-point inequality is known on $Y$ over each field in that list. Apply it with $\eps/2$ and form the finite union of the images of its exceptional loci and their $k$-conjugates. Each image is proper because $f$ is finite. Outside this union and $D$, choose a lift $Q$ of $P$. Functoriality gives
\[
 m_{D,S}(P)+h_{K_X}(P)
 =m_{D_Y,S_L}(Q)+h_{K_Y}(Q)+O(1),\qquad
 h_{A_Y}(Q)=h_A(P)+O(1).
\]
The maximum of the finitely many constants proves the required inequality downstairs. This does not construct an upstairs inequality, and is not an assertion that field-dependent constants can be maximized over infinitely many fields.

\paragraph{Attempt II: exact ramification cancellation.}
Assume a finite surjective map $f:Y\to X$ between smooth projective varieties satisfies
\begin{equation}\label{r64:cover}
 D_Y=(f^*D)_{\mathrm{red}}\text{ is SNC},\qquad
 R=K_Y-f^*K_X=f^*D-D_Y\ge0.
\end{equation}
These are hypotheses on the cover: all ramification is absorbed by the boundary with the displayed multiplicities. They are not asserted for every finite map. For $P=f(Q)\in X(k)\setminus D$, additivity gives the exact relation modulo bounded functions
\begin{align}
 m_{D_Y,S_L}(Q)+h_{K_Y}(Q)
 &=m_{D,S}(P)-m_{R,S_L}(Q)+h_{K_X}(P)+h_R(Q)+O(1)\notag\\
 &=m_{D,S}(P)+h_{K_X}(P)+N_{R,S_L}(Q)+O(1).
 \label{r64:identity}
\end{align}
The sign of the counting term is essential.

Generalized Chevalley--Weil supplies the established estimate
\begin{equation}\label{r64:cw}
 d_k^{\mathrm{abs}}(Q)-d_k^{\mathrm{abs}}(P)
       \le N_{R,S_L}(Q)+O(1).
\end{equation}
This is Theorem 23.15 in \eref{64}{1}, after dividing by $[k:\Q]$. The difference between its discriminant truncated outside $S$ and the full discriminant is bounded at bounded degree, as explained immediately before Lemma 23.4. Here $[k(Q):k]\le d$ and $d_k^{\mathrm{abs}}(P)=0$.

\begin{notebookproposition}\label{r64:descent}
For the chosen cover, suppose every $\delta>0$ admits one proper closed $Z_Y\subsetneq Y$ and one constant $C_\delta$ such that
\begin{equation}\label{r64:upstairs}
 m_{D_Y,S_L}(Q)+h_{K_Y}(Q)
 \le\delta h_{A_Y}(Q)+d_k^{\mathrm{abs}}(Q)+C_\delta
\end{equation}
for every $Q\notin Z_Y$ of degree at most $d$ with $f(Q)\in X(k)$. Then the target Vojta inequality holds for $(X,D,A,S)$.
\end{notebookproposition}
\begin{proof}
Subtract $N_{R,S_L}(Q)$ in~\eqref{r64:identity}, use~\eqref{r64:upstairs}, and then~\eqref{r64:cw}. The result is
\[
 m_{D,S}(P)+h_{K_X}(P)\le\delta h_A(P)+C'.
\]
Exclude $f(Z_Y)$, its conjugates if necessary, and $D$. Finiteness makes this a proper closed subset. Taking $\delta=\eps/2$, shifting the ample height to be nonnegative, and increasing $C'$ by one yields the required strict inequality.
\end{proof}

The source's Proposition 24.2 already proves a more general functorial statement. Thus the displayed transfer is a rederivation of known machinery in the selected normalization, not a new descent theorem. The proposed research step would be to find useful covers on which~\eqref{r64:upstairs} can actually be proved. Taking the identity cover is merely a reformulation; applying the original conjecture separately over varying fields does not establish the needed common constant and exceptional set.

\paragraph{Stress test: the double cover detects the hidden cost.}
Take $X=Y=\PP^1$, $f(z)=z^2$, and $D=D_Y=\{0,\infty\}$. Then $R=[0]+[\infty]$. For a rational prime $p$, let $P=p$ and $Q=\sqrt p$. With $S$ archimedean and standard heights,
\[
 h_{\mathcal O(1)}(P)=\log p,\quad
 d_{\Q}^{\mathrm{abs}}(Q)=\tfrac12\log p+O(1),\quad
 N_{R,S_L}(Q)=\tfrac12\log p.
\]
The field discriminant is $p$ or $4p$. Its contribution is not $o(h(P))$, and infinitely many quadratic fields occur. It is the counting term, not a small-discriminant assertion, that cancels the cost. This example prevents a proof from silently replacing bounded-degree points by points over a fixed field.

There is also a dependency warning. For positive-dimensional smooth projective general-type $X$, the $D=0$ case implies Bombieri--Lang. Indeed, bigness gives $mK_X\sim A_0+E$ with $A_0$ ample and $E$ effective. Outside $E$, comparison of ample heights gives $h_{K_X}\ge c h_A-O(1)$ for some $c>0$. Choosing $\eps<c$ bounds $h_A$ outside the exceptional set; Northcott then leaves only finitely many such rational points. A universal closing argument here would therefore also resolve the earlier Problem 55, rather than only a normalization issue.

\paragraph{Outcome.}
\textbf{Outcome: Conditional reduction.}

The attempt isolates a uniform bounded-degree inequality on a chosen ramified cover that suffices for the original rational-point assertion. Checking the source removes one apparent missing estimate: discriminant control is already supplied by generalized Chevalley--Weil. No useful family of covers with a proved upstairs inequality has been produced. The transfer itself is known, its specialization and stress test are explicit here, and no full resolution or new general height theorem is claimed.

% END RESEARCH INSERTION
\subsection{Sources}
\begin{itemize}\raggedright
\item[\textnormal{[S1]}]\hypertarget{source:64:1}{} Yasufuku, GCD inequalities inspired by Vojta's conjecture, Conjecture 8 (2024).
\url{https://link.springer.com/article/10.1007/s00605-024-02018-1}.
{\small\textit{Catalogue formulation source.}}
\item[\textnormal{[S2]}]\hypertarget{source:64:2}{} Towards Lang--Vojta via Degeneration.
\url{https://arxiv.org/abs/2602.06956}.
{\small\textit{Catalogue research/status reference; not a proof endorsement.}}
\item[\textnormal{[S3]}]\hypertarget{source:64:3}{} Lang-Vojta conjecture over function fields for very general log projective spaces.
\url{https://arxiv.org/abs/2606.14074}.
{\small\textit{Catalogue research/status reference; not a proof endorsement.}}

% BEGIN ADDED RESEARCH REFERENCES
\item[\textnormal{[E1]}]\hypertarget{extra:64:1}{} Paul Vojta, \emph{Diophantine Approximation and Nevanlinna Theory}, expanded CIME lecture notes, author-hosted version; Sections 12, 14, and 23--24, especially Definition 14.1, Theorem 23.15, Proposition 24.2, and the discussion preceding Lemma 23.4.
\url{https://math.berkeley.edu/~vojta/cime/cime.pdf}.
{\small\textit{Generalized Chevalley--Weil, bounded-degree discriminants, and known functoriality of the height inequality. Consulted 27 September 2026.}}
\item[\textnormal{[E2]}]\hypertarget{extra:64:2}{} Ryan C. Chen, Natalia Garcia-Fritz, Siddharth Mathur, and Hector Pasten, \emph{Towards Lang--Vojta via Degeneration}, arXiv:2602.06956 (6 February 2026).
\url{https://arxiv.org/abs/2602.06956}.
{\small\textit{Recent results for specified integral-point problems; not used as a universal rational-point height bound. Consulted 27 September 2026.}}
\item[\textnormal{[E3]}]\hypertarget{extra:64:3}{} Takeshi Abe, \emph{Lang-Vojta conjecture over function fields for very general log projective spaces}, arXiv:2606.14074 (12 June 2026).
\url{https://arxiv.org/abs/2606.14074}.
{\small\textit{Function-field setting, kept separate from the number-field target. Consulted 27 September 2026.}}
% END ADDED RESEARCH REFERENCES
\end{itemize}

\end{document}
